% =============================================================================
% Methodology
% =============================================================================
\section{Methodology}\label{sec:methodology}

\subsection{Jacobian computation}

The Jacobian~\eqref{eq:jacobian} is computed by central finite differences:
\begin{equation}
  J_{ij}
  \approx \frac{f_i(\boldsymbol{\theta} + \delta_j \mathbf{e}_j)
              - f_i(\boldsymbol{\theta} - \delta_j \mathbf{e}_j)}
               {2 \delta_j},
\end{equation}
where $\delta_j = \epsilon \lvert \theta_j \rvert$ with $\epsilon = 10^{-6}$.
Convergence is verified by comparing step sizes $\epsilon = 10^{-5}$ and
$10^{-7}$; agreement to four significant figures confirms that the
perturbation is within the smooth regime.

\subsection{Condition number}

The condition number of the Jacobian in the 2-norm,
\begin{equation}
  \kappa(\mathbf{J}) = \frac{\sigma_\text{max}}{\sigma_\text{min}},
\end{equation}
where $\sigma_i$ are the singular values of~$\mathbf{J}$, quantifies the
worst-case amplification of measurement noise into parameter uncertainty.
A value $\kappa \sim 1$ indicates ideal identifiability;
$\kappa \gtrsim 10^6$ indicates practical non-identifiability.

\subsection{Fisher information and Cram\'er--Rao bounds}

Assuming additive Gaussian noise with covariance
$\boldsymbol{\Sigma} = \operatorname{diag}(\sigma_1^2, \dots, \sigma_N^2)$,
where $\sigma_i = \epsilon_\text{noise} \, f_i$, the Fisher information
matrix is
\begin{equation}\label{eq:fisher}
  \mathbf{F}
  = \mathbf{J}^\top \boldsymbol{\Sigma}^{-1} \mathbf{J}.
\end{equation}
The Cram\'er--Rao lower bound on the variance of any unbiased estimator of
$\theta_j$ is
\begin{equation}
  \operatorname{Var}(\hat\theta_j) \geq [\mathbf{F}^{-1}]_{jj}.
\end{equation}
When $\mathbf{F}$ is singular (as for the sphere model), no finite-variance
unbiased estimator exists, confirming non-identifiability.

\subsection{Newton--Raphson inversion}

We recover parameters by minimising the squared residual
\begin{equation}
  \mathcal{C}(\boldsymbol{\theta})
  = \tfrac{1}{2}
    \lVert \mathbf{f}(\boldsymbol{\theta}) - \mathbf{f}^\text{obs} \rVert^2
\end{equation}
using a trust-region reflective algorithm (Levenberg--Marquardt variant) with
positivity constraints on all parameters.  The initial guess is perturbed by
$\pm 10\%$ from the true values to test convergence basins.

\subsection{Parameter-space condition maps}

To assess how identifiability varies across physiologically relevant
parameter ranges, we sweep $\kappa(\mathbf{J})$ over a three-dimensional
grid: $a \in [\SI{0.14}{m}, \SI{0.22}{m}]$,
$c \in [\SI{0.08}{m}, \SI{0.16}{m}]$,
$E \in [\SI{0.05}{MPa}, \SI{0.5}{MPa}]$, holding all other parameters at
their canonical values.
